\section{Conclusiones}

Para encontrar la solución del problema, es importante destacar que en programación dinámica es necesario realizar un análisis previo del problema, siendo este el mayor problema del mismo. Teniendo en cuenta los algoritmos de Greedy que se tiene que elegir la decisión codiciosa que nos puede dar la mejor solución local en el estado correspondiente, en programación dinámica buscamos encontrar la mejor solución del estado actual definiendo mis variables de estados, de esta forma puedo guardar o memorizar los valores previos para poder ahorrar re calcular el mismo estado varias veces, además la programación dinámica nos permite explorar diferentes caminos para poder encontrar la mejor solución.

Es importante que reconstruir la solución tiene su complejidad que tenemos que asegurar que el tiempo que demore a lo sumo sea igual que encontrar la solución, en el caso que sea peor, seria mejor reconstruir el algoritmo porque encontrar la solución debería ser el algoritmo con mayor complejidad del problema.

Un punto extra es la medición de los tiempos del algoritmo, esto nos permite chequear que la teoría se corroborar con la practica realizada, usando métodos matemáticos que nos permiten buscar una función que mejor aproxima los puntos medidos desde un código aparte.